%==============================================================================
%  17-case-studies.tex — Benchmark Case Studies: Success & Failure
%==============================================================================
\strategyPart{17 \textbar\ Case Studies}{Success and Failure Stories with Measured Results}

The strongest evidence for the problem-first philosophy comes from cases with
measured outcomes — in both directions. This part compares documented
successes and failures.

\section{Success: Morgan Stanley — knowledge reinvention}

Morgan Stanley identified an expensive organizational problem: \~{}50,000
financial advisors could not search the firm's research library fast enough
to answer client questions in real time. In 2023 it launched \textit{AI @
Morgan Stanley Assistant}, built on OpenAI's GPT-4, giving advisors a
conversational interface to internal research [Morgan Stanley 2023; Business
Insider 2023]. A Level-1 (Assist) deployment with a precise problem
definition and a measurable outcome.

\section{Success: Coca-Cola — creative and supply-chain reinvention}

Coca-Cola's \textit{Create Real Magic} generative-AI platform — built with
OpenAI and Bain — attacks a content-production bottleneck across a global
brand portfolio [Forbes 2023], while predictive and optimization AI support
demand forecasting and supply-chain planning. In 2024 the company committed
over \$1~billion to a five-year cloud and generative-AI partnership with
Microsoft [Reuters 2024].

\section{Success: Klarna — customer-service automation}

Klarna's AI assistant, launched in 2024, handled roughly two-thirds of
customer-service chats in its first month — the equivalent workload of about
\textbf{700 full-time agents} — while maintaining comparable customer
satisfaction and resolution times [Klarna 2024; press coverage]. The case
demonstrates Level-2 (Automate) value with a precisely scoped workflow.

\section{Failure: Air Canada — an unreliable assistant}

In 2024 a Canadian tribunal held Air Canada liable for incorrect information
provided by its AI chatbot about bereavement fares — because the company was
responsible for the chatbot's outputs [tribunal ruling, 2024]. The failure was
not the model; it was governance: no policy defined what the assistant could
claim and who owned its outputs.

\section{Failure: IBM Watson Health — technology-first at scale}

IBM invested heavily in Watson for healthcare, but the system struggled to
move from demonstration to clinical reality; the business was eventually sold
off in 2022 after years of losses [published analyses]. The case is the
canonical example of a technology-first strategy — capability acquired before
a defined, governable problem.

\section{Failure: Zillow Offers — scaling an unvalidated model}

Zillow's algorithmic home-buying business scaled a predictive model whose
error rate exceeded the trading margin, producing an \textbf{\$881~million
write-down} and the unit's closure in 2021 [company filings; press]. Scaling a
use case without validating risk and reliability is the eighth common failure
pattern — here, with a billion-dollar price tag.

\section{Synthesis: problems beat tools}

Across the research base, organizations that frame AI around business problems
and KPIs reach production and scale; organizations that frame AI around
technology stall at proof of concept — the population captured by Gartner's
30\% abandonment prediction [Gartner 2023]. Roughly three-quarters of
generative-AI value sits in a handful of functions where the problems are
biggest [McKinsey 2023]. \textbf{The organization that starts with the problem
wins; the organization that starts with the tool waits — or loses.}

